\section{When can synthetic pre-pre-training help?}
\label{sec:transfer}

\subsection{The effect of resetting the token table}

Write $f(B)=L_{\rm lang}(B,E_0)$ for language loss immediately after
transfer. A synthetic stage may train both $B$ and its token table, but
the transfer map keeps only the backbone change. Its differential is
\begin{equation}
 D\mathcal T\,(\delta B,\delta E)=(\delta B,0).
 \label{eq:projection}
\end{equation}
The synthetic table may even have a different vocabulary size; the
backbone comparison is still well-defined. Synthetic head learning can
affect later backbone updates \emph{during} pre-pre-training, but its
terminal displacement is not itself carried across the reset.

Let $B_G=B_0+\delta_G$ be any synthetic-trained backbone. Assume $f$ has
an $L$-Lipschitz gradient on the segment joining $B_0$ and $B_G$.
Taylor's integral remainder gives
\begin{equation}
 \begin{aligned}
 f(B_G)-f(B_0)&=g^\top\delta_G+R_G,\qquad g=\nabla f(B_0),\\
 R_G&=\int_0^1(1-t)\delta_G^\top
    \nabla^2 f(B_0+t\delta_G)\delta_G\,dt.
 \end{aligned}
\end{equation}
Consequently
\begin{equation}
 \left|f(B_G)-f(B_0)-g^\top\delta_G\right|
 \leq\frac L2\norm{\delta_G}^2.
 \label{eq:displacement}
\end{equation}
This is local and requires no claim that a transformer loss is globally
convex.

\begin{proposition}[Sufficient backbone-transfer condition]
\label{prop:alignment}
Under the preceding smoothness assumption,
\begin{equation}
 -g^\top\delta_G>\frac L2\norm{\delta_G}^2
 \quad\Longrightarrow\quad f(B_G)<f(B_0).
 \label{eq:alignment}
\end{equation}
\end{proposition}
\begin{proof}
Use the upper side of \eqref{eq:displacement}. The right side of
$f(B_G)-f(B_0)\leq g^\top\delta_G+(L/2)\norm{\delta_G}^2$
is strictly negative under \eqref{eq:alignment}.
\end{proof}
For $\delta_G\ne0$, $g\ne0$, define
$a_G=-g^\top\delta_G/(\norm g\norm{\delta_G})$.
The sufficient condition is
$a_G>L\norm{\delta_G}/(2\norm g)$.
Positive direction alignment can be destroyed by too large a step.
More synthetically learned structure is not automatically a better
initialization.

For one synthetic gradient step, $\delta_G=-\eta g_G^{\rm syn}$, so
\begin{equation}
 g^\top g_G^{\rm syn}>
       \frac{L\eta}{2}\norm{g_G^{\rm syn}}^2
 \label{eq:singlealignment}
\end{equation}
is sufficient. The language and synthetic backbone gradients are
evaluated at the shared backbone, with their respective initial token
tables. For Adam or another optimizer, replace $g_G^{\rm syn}$ by its
actual update direction $v_G$. For a fixed diagonal preconditioner
$A\succ0$, the relevant alignment is $g^\top A g_G^{\rm syn}$;
ordinary Euclidean gradient cosine need not have the same sign.
Weight decay, clipping and momentum must also be included in $v_G$.
For a multi-step synthetic stage, \eqref{eq:displacement} applies to its
total backbone displacement directly; it is not necessary to pretend
Adam performed a single gradient step.

For grid and graph one-step displacements, two-sided remainder bounds
give the stronger comparison
\begin{align}
 f(B_G)-f(B_{\rm grid})
 &\leq-\eta\,g^\top
           (g_G^{\rm syn}-g_{\rm grid}^{\rm syn})\\
 &\quad+\frac{L\eta^2}{2}
       \left(\norm{g_G^{\rm syn}}^2+
             \norm{g_{\rm grid}^{\rm syn}}^2\right).
 \label{eq:graphvsgrid}
\end{align}
A larger first-order graph alignment must exceed the displayed
remainder to certify an advantage by this sufficient bound.
Subtracting two separate upper bounds would not justify
\eqref{eq:graphvsgrid}; its derivation uses an upper remainder for graph
and a lower remainder for grid.

\subsection{An explicit topology-to-transfer calculation}

The preceding criterion leaves the synthetic gradient unspecified.
Here is a model in which it can be derived from the graph exactly.
It is a squared-loss linear prediction surrogate, not a claim that
binary NCA token prediction is linear.
Let $x$ have mean zero and covariance $I_n$; independent fair signed
bits are sufficient. Define a degree-normalized graph target
\[
 A_G=\frac1d\sum_{r=1}^dP_r,\qquad y_G=A_Gx.
\]
A trainable matrix $M$ predicts $Mx$, with population synthetic loss
\begin{equation}
 L_G^{\rm lin}(M)=\frac12\E\norm{(M-A_G)x}_2^2
                 =\frac12\norm{M-A_G}_F^2.
 \label{eq:linearsynthetic}
\end{equation}
The equality follows from
$\E[x^\top R^\top R x]={\rm tr}(R^\top R\E[xx^\top])$.
Starting at $M=0$, $K$ synthetic gradient updates with
$0<\eta<1$ give the exact formula
\begin{equation}
 M_G=\gamma A_G,\qquad
 \gamma=1-(1-\eta)^K.
 \label{eq:linearpretrain}
\end{equation}
For a target task $y_T=Cx$, define
$L_T^{\rm lin}(M)=\tfrac12\norm{M-C}_F^2$. Its change from scratch is
\begin{equation}
 L_T^{\rm lin}(M_G)-L_T^{\rm lin}(0)
 =\frac{\gamma^2}{2}\norm{A_G}_F^2
       -\gamma\langle A_G,C\rangle_F.
 \label{eq:graphcovtransfer}
\end{equation}
This supplies a concrete interpretation: topology enters through the
cross-covariance $\E[y_Gx^\top]=A_G$, and useful transfer depends on
its overlap with the target operator $C$.

When every row has $d$ distinct neighbors, each row of $A_G$ has $d$
entries equal to $1/d$, so $\norm{A_G}_F^2=n/d$.
For a degree-matched grid operator $A_0$, these norm terms cancel:
\begin{equation}
 L_T^{\rm lin}(\gamma A_G)-L_T^{\rm lin}(\gamma A_0)
 =-\gamma\langle A_G-A_0,C\rangle_F.
 \label{eq:exactgraphtarget}
\end{equation}
If $C=A_G$, this becomes
$-\tfrac{\gamma}{2}\norm{A_G-A_0}_F^2$.
If $C=A_0$, its sign reverses. Degree matching alone therefore cannot
choose the favorable target.

For a fully numerical construction, take $n=6,d=2$.
Let $A_0$ average cyclic offsets $\pm1$ modulo six, and let $A_G$
average offsets $\pm2$. Both are valid two-regular typed graphs;
the latter consists of two triangles. Their edge supports are disjoint:
\[
 \norm{A_0}_F^2=\norm{A_G}_F^2=3,\qquad
 \langle A_0,A_G\rangle_F=0.
 \]
Use one synthetic step with $\eta=1/2$, so $\gamma=1/2$, and set
$C=A_G$. Target losses are then
\begin{align}
 L_T^{\rm lin}(0)&=\frac32=1.5,\\
 L_T^{\rm lin}(\tfrac12 A_G)&=\frac12\cdot\frac14\cdot3
                              =0.375,\\
 L_T^{\rm lin}(\tfrac12 A_0)&=\frac12
                              \left(\frac14\cdot3+3\right)
                              =1.875.
\end{align}
The graph-matched synthetic task helps and the grid-like cycle task
hurts in this constructed target. Setting $C=A_0$ reverses the
roles. A disconnected graph can thus be useful for a matching target
even though it has poor global diffusion mixing.
This proves a conditional mechanism in a fully specified surrogate,
not an improvement for natural language or the actual nonlinear NCA.

\subsection{A completely worked gradient-alignment example}

Consider the constructed loss on $\R^2$
\begin{equation}
 f(b)=\frac12(b-b_*)^\top(4I)(b-b_*),
 \quad b_*=(-1/2,-1/4),\quad b_0=(0,0).
\end{equation}
Then $g=(2,1)$, $L=4$, and $f(b_0)=5/8=0.625$.
Take $\eta=0.1$, graph synthetic gradient $(1,1)$, and grid gradient
$(-1,1)$. These vectors are specified examples, not measured project
gradients. The exact quadratic changes are
\begin{align}
 \Delta f_G&=-0.1(2+1)+
       \frac42(0.1)^2(1^2+1^2)=-0.26,\\
 \Delta f_{\rm grid}&=-0.1(-2+1)+
       \frac42(0.1)^2((-1)^2+1^2)=+0.14.
\end{align}
Thus $f_G=0.365$ and $f_{\rm grid}=0.765$.
Equal synthetic gradient norms and update counts do not imply equal
transfer; their alignment with the target matters. Interchanging the
two specified synthetic gradients interchanges the conclusion. Graph
topology has not supplied either favorable gradient by theorem.

\subsection{Exact convergence in a local quadratic model}

Now let $\theta$ denote all language parameters after transfer, including
the common token table. Assume an exactly quadratic objective on a
chosen parameter subspace:
\begin{equation}
 f(\theta)-f_*=\frac12(\theta-\theta_*)^\top
                   H(\theta-\theta_*),\qquad
 0<\mu I\preceq H\preceq LI.
 \label{eq:quadratic}
\end{equation}
This is an analytic model, not an assertion about the global loss of a
transformer. Under full-batch gradient descent with $0<\alpha<2/L$,
$e_n=\theta_n-\theta_*$ obeys
\begin{equation}
 e_n=(I-\alpha H)^n e_0.
\end{equation}
Using orthonormal eigenvectors $u_j$ of $H$, let
$c_j=u_j^\top e_0$ and $H u_j=\lambda_j u_j$. Then
\begin{equation}
 \boxed{\ f_n-f_*=
   \frac12\sum_j \lambda_j(1-\alpha\lambda_j)^{2n}c_j^2.\ }
 \label{eq:modeformula}
\end{equation}
This follows by substituting the eigendecomposition into the quadratic
form. No fitted rate constant or simulated trajectory is needed.

In particular, if transferred coefficients satisfy
$(c_j^G)^2\leq(c_j^0)^2$ for every $j$, then transferred loss is no worse
at every $n$. A smaller total initial loss is a weaker condition and
does not suffice: a transfer can increase error along slowly decaying
modes while reducing fast-mode error.

\subsection{Counterexample: initially better, subsequently worse}

Take $H=\diag(1,100)$, $\alpha=0.01$, scratch error $e_0=(1,0.2)$ and
transferred error $e_G=(2,0)$. Initially,
\begin{equation}
 f_0-f_*=\frac12(1+100\cdot0.2^2)=2.5,\qquad
 f_G-f_*=\frac12(2^2)=2.
\end{equation}
The fast factor $1-\alpha100$ is zero. After one common language step,
\begin{equation}
 f_1-f_*=\frac12(0.99)^2=0.49005,\qquad
 f_{G,1}-f_*=\frac12(2\cdot0.99)^2=1.9602.
\end{equation}
The initial ranking reverses exactly. These coordinates can be
backbone directions, so the example is compatible with a common,
unchanged token-table initialization.

There is also no universal improvement theorem based only on graph
structure: choose a target quadratic minimized at the scratch
initialization and any synthetic displacement away from it. Scratch
has zero excess loss and stays optimal, while the displaced model has
positive excess loss for every finite non-annihilating gradient step.
Additional assumptions linking synthetic updates to the target are
therefore logically necessary.

